% ============================================================
% TechRxiv / IEEE Preprint
% Obsidian AI Assistant
% ============================================================
\documentclass[10pt,a4paper]{article}

% --- Packages ---
\usepackage[margin=1in]{geometry}
\usepackage{hyperref}
\usepackage{graphicx}
\usepackage{float}
\usepackage{booktabs}
\usepackage{amsmath}
\usepackage{amssymb}
\usepackage{xcolor}
\usepackage{listings}
\usepackage{url}

\hypersetup{
  colorlinks=true,
  linkcolor=blue,
  urlcolor=blue,
  citecolor=blue,
}

% --- Title ---
\title{Obsidian AI Assistant: An LLM-Powered Knowledge Management Plugin\\with CJK-Optimized Semantic Search and Autonomous Tool-Calling}

\author{
  Yiqi Cai\textsuperscript{1}\thanks{Corresponding author: yiqi.cai@example.com}\\[4pt]
  \textsuperscript{1} School of Computer Science, South-Central Minzu University, Wuhan, China
}

\date{July 2026}

\begin{document}
\maketitle

% --- Abstract ---
\begin{abstract}
Obsidian AI Assistant is an open-source plugin for Obsidian---a popular local-first knowledge management tool with over one million users. The plugin transforms Obsidian into an AI-augmented workspace by enabling large language models (LLMs) to autonomously search, create, modify, and organize notes within the user's vault. It supports multiple Chinese LLM providers (DeepSeek, Qwen, GLM-4) and local models via Ollama, addresses CJK (Chinese/Japanese/Korean) text processing challenges absent in existing solutions, and provides a robust agentic tool-calling framework with multi-layered error recovery. Benchmark evaluation demonstrates that our CJK n-gram tokenizer achieves 83.3\% Precision@5 (vs.\ 25.5\% substring baseline), while the Canvas JSON normalizer and tool-call parser achieve 100\% robustness across 22 adversarial test cases. The plugin is distributed via GitHub Releases under the MIT license and archived at Zenodo (DOI: 10.5281/zenodo.21442160).
\end{abstract}

\noindent\textbf{Keywords:} Obsidian plugin, LLM, knowledge management, CJK tokenization, agentic tool-calling, RAG, Canvas, DeepSeek, Qwen, GLM

% ============================================================
\section{Introduction}

Knowledge management tools have become essential for researchers, developers, and writers who accumulate large volumes of interconnected notes. Obsidian---a local-first, Markdown-based knowledge base with bidirectional linking---has attracted over one million users since its launch~\cite{obsidian}. Its plugin ecosystem enables community-developed extensions, with AI-assisted plugins representing one of the most active categories.

Existing AI plugins for Obsidian, such as Copilot~\cite{obsidian_copilot}, predominantly target English-speaking users and rely on OpenAI's API. This creates three barriers for Chinese-speaking knowledge workers: (1) network inaccessibility to OpenAI services, (2) privacy concerns with uploading personal notes to foreign servers, and (3) inadequate CJK text search due to whitespace-based tokenization. Furthermore, current plugins lack autonomous tool-use capabilities---they function as passive Q\&A interfaces rather than active knowledge management agents capable of modifying the vault.

Obsidian AI Assistant addresses these gaps by providing: (a) native support for Chinese LLM providers with a unified adapter pattern, (b) client-side CJK n-gram tokenization and hybrid TF-IDF/embedding search, and (c) an agentic loop enabling the AI to autonomously execute nine vault operations.

\begin{table}[H]
\centering
\caption{Feature comparison between Obsidian Copilot and the proposed plugin.}
\label{tab:comparison}
\begin{tabular}{lcc}
\toprule
\textbf{Feature} & \textbf{Obsidian Copilot} & \textbf{Ours} \\
\midrule
Chinese LLM providers (DeepSeek/Qwen/GLM/Ollama) & \textcolor{red}{$\times$} & \textcolor{green!60!black}{\checkmark} \\
Autonomous vault operations (Agentic Loop) & \textcolor{red}{$\times$} & \textcolor{green!60!black}{\checkmark} \\
CJK-optimized full-text search & \textcolor{red}{$\times$} & \textcolor{green!60!black}{\checkmark} \\
Knowledge graph generation (Canvas) & \textcolor{red}{$\times$} & \textcolor{green!60!black}{\checkmark} \\
Client-side PII sanitization & \textcolor{red}{$\times$} & \textcolor{green!60!black}{\checkmark} \\
Local ONNX embedding (offline mode) & \textcolor{red}{$\times$} & \textcolor{green!60!black}{\checkmark} \\
Multi-layer error recovery for AI output & \textcolor{red}{$\times$} & \textcolor{green!60!black}{\checkmark} \\
Built-in Zettelkasten/PARA methodology in prompts & \textcolor{red}{$\times$} & \textcolor{green!60!black}{\checkmark} \\
\bottomrule
\end{tabular}
\end{table}

\section{System Design}

\subsection{Multi-Provider LLM Router}

The \texttt{ChatModelManager} implements a unified adapter for four LLM providers (DeepSeek~\cite{deepseek}, Qwen, GLM-4, Ollama) using an OpenAI-compatible protocol. Provider-specific differences (base URL, API key, model name) are abstracted behind a single \texttt{chat()} interface. A separate vision model pipeline routes image, PDF, and Word files to multimodal models (Qwen-VL, GLM-4V) for OCR and content extraction.

\subsection{Agentic Tool-Calling Loop}

The core innovation is an autonomous agent loop (up to 20 rounds) in which the AI receives user input, optionally calls vault manipulation tools, processes results, and iterates until producing a final answer. Unlike OpenAI's native Function Calling, which is inconsistently supported across Chinese LLMs, the plugin uses a custom \texttt{<tool\_call>} XML protocol parsed via regular expressions, with a fallback parser for GLM-4's malformed output patterns.

Nine built-in tools are registered in a singleton \texttt{ToolRegistry}: \texttt{listNotes}, \texttt{searchVault}, \texttt{readNote}, \texttt{createNote}, \texttt{modifyNote}, \texttt{appendNote}, \texttt{getFileTree}, \texttt{getTags}, and \texttt{saveCanvas}. A safety net automatically detects when the AI has produced note content but forgotten to call \texttt{createNote}, preventing data loss.

\subsection{CJK-Optimized Semantic Search}

The plugin implements a lightweight TF-IDF index using a custom CJK n-gram tokenizer. Chinese text contains no whitespace delimiters, rendering standard word-boundary tokenizers ineffective. Our tokenizer performs bigram and trigram extraction on CJK runs while falling back to word-boundary splitting for non-CJK segments. Title tokens receive a 3$\times$ weighting boost, and only the first 5,000 characters of each note are indexed for performance.

An optional hybrid search mode combines TF-IDF scores with embedding-based semantic retrieval (via Qwen/GLM Embedding APIs or ONNX local models) using weighted fusion:
\begin{equation}
\text{score} = 0.3 \times \text{TF-IDF}_{\text{norm}} + 0.7 \times \text{Embedding}_{\text{norm}}
\end{equation}

\subsection{Canvas JSON Error Recovery}

AI-generated Canvas JSON is notoriously unreliable---different models produce inconsistent node types, interleave edges within node arrays, reference nodes by text labels instead of IDs, and output duplicate edges. The \texttt{normalizeCanvasJSON} function implements five error-recovery strategies: (1) multi-candidate bracket-counting extraction from arbitrary text, (2) edge-node separation via \texttt{fromNode} property detection, (3) non-standard type normalization to \texttt{"text"}, (4) label-to-ID mapping for edge references, and (5) automatic grid layout fallback for missing coordinates.

\subsection{Privacy and Knowledge Management}

Client-side PII sanitization (phone numbers, ID cards, emails, IP addresses) runs locally before any API transmission. A memory system with LRU eviction caches conversation context. The system prompt encodes Zettelkasten atomic note principles~\cite{zettelkasten}, PARA folder organization methodology, and MOC (Maps of Content) indexing strategies.

\section{Design Decisions}

The plugin embodies a series of non-trivial design decisions that distinguish it from AI-generated boilerplate, summarized in Table~\ref{tab:decisions}.

\begin{table}[H]
\centering
\caption{Key Design Decisions and Their Rationale}
\label{tab:decisions}
\begin{tabular}{p{3.2cm}p{6.8cm}}
\toprule
\textbf{Decision} & \textbf{Rationale} \\
\midrule
Custom \texttt{<tool\_call>} XML over Function Calling & GLM-4 Function Calling is unstable; regex parsing is provider-agnostic \\
GLM-4 fallback parser & Malformed output pattern observed in production; fallback prevents silent failures \\
Multi-layer Canvas JSON recovery & AI JSON output is unreliable; 5 error types handled across all supported models \\
CJK n-gram tokenizer (no jieba) & Minimizes bundle size; dependency-free and adequate for note-scale text \\
Auto-save safety net & AI occasionally forgets tool calls; content pattern detection prevents data loss \\
Stream progress cards & Real-time feedback for long generations; unique among Obsidian plugins \\
Client-side PII sanitization & Privacy-by-design; prevents sensitive data from reaching third-party APIs \\
Zettelkasten + PARA + MOC prompts & Encodes established knowledge management methodology into AI behavior \\
LRU memory eviction (80\% threshold) & Prevents infinite growth; hysteresis avoids thrashing \\
3-level search degradation & Graceful fallback: API $\rightarrow$ ONNX $\rightarrow$ TF-IDF \\
\bottomrule
\end{tabular}
\end{table}

\section{Benchmark Evaluation}

We conduct four controlled experiments to quantify the plugin's robustness.

\subsection{Tokenizer Precision}

Evaluated against 15 test cases covering short queries, Chinese-English mixed input, punctuation, and edge cases. N-gram tokenization achieves 100\% recall, while whitespace-based tokenization fails entirely on Chinese text.

\subsection{Canvas JSON Robustness}

Ten test cases simulate common LLM output errors: edges in node arrays, non-standard types, missing coordinates, label-based references, and duplicates. The normalizer correctly handles all error types (100\% pass).

\subsection{Tool Call Parser Robustness}

Twelve cases evaluate parsing against standard \texttt{<tool\_call>} blocks, GLM-4 malformed output, nested JSON, escaped quotes, and corrupted input. The dual-path parser achieves 100\% accuracy.

\subsection{Search Precision}

Using a simulated vault of 30 notes and 49 queries across four categories (short, long, mixed Chinese-English, edge cases), we compare CJK n-gram TF-IDF against two baselines. Results in Table~\ref{tab:search}.

\begin{table}[H]
\centering
\caption{Search Precision@5 and Recall@5 Comparison}
\label{tab:search}
\begin{tabular}{lcc}
\toprule
\textbf{Method} & \textbf{Precision@5} & \textbf{Recall@5} \\
\midrule
CJK n-gram TF-IDF (ours) & 83.3\% & 89.5\% \\
Substring match (baseline) & 25.5\% & 29.6\% \\
Dictionary match (baseline) & 72.1\% & 78.2\% \\
\bottomrule
\end{tabular}
\end{table}

\section{Availability}

The plugin is available on GitHub at \url{https://github.com/Qi-hub-dot/obsidian-multi-llm-ai-agent} under the MIT license. The repository includes 131 unit tests (Jest 30), architecture documentation, and a contribution guide. Version 2.0.7 is archived at \url{https://doi.org/10.5281/zenodo.21442160}.

\section{Conclusion}

Obsidian AI Assistant fills a gap in the Obsidian ecosystem by providing first-class support for Chinese LLM providers, CJK-optimized semantic search, and autonomous agentic tool-calling. The multi-layer error recovery mechanism ensures robustness against real-world LLM output variability---a practical concern rarely addressed in plugin development. Benchmark results demonstrate significant improvements over baseline approaches, achieving 3.3$\times$ search precision improvement over substring matching. Future work includes expanding provider support, implementing local embedding inference via WebGPU, and integrating with additional knowledge management paradigms.

\section{Acknowledgments}

This project was developed independently. The architecture design drew inspiration from Obsidian Copilot~\cite{obsidian_copilot}. Benchmark evaluation methodology follows best practices from the information retrieval community~\cite{joss}.

% --- References ---
\bibliographystyle{IEEEtran}
\bibliography{../paper}

\end{document}
